\documentclass[11pt]{article}

% ============================================================
%  PLANTILLA MINIMALISTA (adaptada de la plantilla del alumno)
%  - Texto gris oscuro (no negro puro)
%  - Secciones centradas, en mayúsculas, con línea divisoria
%  - Subsecciones discretas
%  - Márgenes estrechos
%  - Footer con nombre (derecha) y número de página (centro)
%  - Header vacío
%  NOTA: se removieron \usepackage{lmodern} y \usepackage{soul}
%  porque no están disponibles en este entorno de compilación y
%  no son necesarios para este documento (no usa \hl, y la fuente
%  Computer Modern por defecto ya cubre lo que se necesita).
% ============================================================

\usepackage[utf8]{inputenc}
\usepackage[T1]{fontenc}
% NOTA: se removió \usepackage[spanish]{babel} porque los patrones de
% hifenación en español no están instalados en este entorno de
% compilación (texlive-lang-spanish) y no hay conexión para instalarlos.
% El documento no depende de macros específicas de babel, así que esto
% no afecta el contenido; solo puede haber una hifenación algo menos
% prolija en palabras largas.
\usepackage{geometry}
\usepackage{amsmath}
\usepackage{amssymb}
\usepackage{enumitem}
\usepackage{tabularx}
\usepackage{booktabs}
\renewcommand{\familydefault}{\sfdefault}
\usepackage{titlesec}
\usepackage{fancyhdr}
\usepackage{ragged2e}
\usepackage[table]{xcolor}
\usepackage{array}

\geometry{
    top=2cm,
    bottom=2.2cm,
    left=2cm,
    right=2cm
}

\definecolor{textgray}{RGB}{60, 60, 60}
\definecolor{rulegray}{RGB}{150, 150, 150}
\color{textgray}

\newcommand{\miNombre}{Villar Pedro}

\pagestyle{fancy}
\fancyhf{}
\renewcommand{\headrulewidth}{0pt}
\renewcommand{\footrulewidth}{0pt}
\fancyfoot[C]{\small\color{textgray}\thepage}
\fancyfoot[R]{\footnotesize\color{textgray}\miNombre}

\titleformat{\section}
  {\centering\color{textgray}\fontsize{13}{15}\selectfont\bfseries}
  {}
  {0pt}
  {}
  [\vspace{2pt}\color{rulegray}\titlerule\vspace{4pt}]

\titlespacing*{\section}{0pt}{18pt}{6pt}

\let\oldsection\section
\renewcommand{\section}[1]{\oldsection{\MakeUppercase{#1}}}

\titleformat{\subsection}
  {\color{textgray}\normalsize\bfseries}
  {}
  {0pt}
  {}

\titlespacing*{\subsection}{0pt}{12pt}{4pt}

\newcommand{\ok}{\checkmark}

% Como no está disponible babel en español en este entorno, se renombra
% manualmente el prefijo de las tablas (por defecto en inglés "Table").
\renewcommand{\tablename}{Tabla}

\begin{document}

\begin{center}
{\fontsize{16}{20}\selectfont\bfseries\color{textgray} ÁLGEBRA EN FEC (II) - GUÍA DE EJERCICIOS}
\end{center}

\vspace{4pt}

\section{Solución}

\subsection*{Ejercicio 1: Polinomios sobre $GF(2)$}

Dadas las secuencias $a=1101010$ y $b=01101101$ (el bit más a la izquierda es el coeficiente
de $X^0$:

\textbf{a)} Leyendo bit a bit,
$$a(X) = 1 + X + X^3 + X^5, \qquad b(X) = X + X^2 + X^4 + X^5 + X^7.$$

\textbf{b)} La suma es coeficiente a coeficiente módulo:
$$a(X) + b(X) = (1+X+X^3+X^5) + (X+X^2+X^4+X^5+X^7) = 1 + X^2 + X^3 + X^4 + X^7,$$
Como secuencia de bits:
$$a+b \;=\; \texttt{10111001}.$$

\textbf{c)} 
$$a(X)\cdot b(X) = X^{12} + X^9 + X^7 + X^5 + X^3 + X.$$

\textbf{d)}
$$q(X) = b(X) = X^7+X^5+X^4+X^2+X, \qquad r(X) = 0. \quad $$


\subsection*{Ejercicio 2 -- División de polinomios}
 
Sean $f(X) = 1+X^2+X^3+X^5+X^6$ y $g(X)=1+X+X^3$.
 
\textbf{a)} 
\bigskip
\renewcommand{\arraystretch}{1.3}
 
\textbf{Paso 1.} El término líder de $f$ es $X^6$; como el de $g$ es $X^3$, el cociente
suma el término $X^{6-3}=X^3$. Se resta $X^3\cdot g(X) = X^6+X^4+X^3$:
$$
\begin{array}{r|c c c c c c c}
 & X^6 & X^5 & X^4 & X^3 & X^2 & X^1 & X^0 \\
\hline
\text{resto}_0=f(X): & 1 & 1 & 0 & 1 & 1 & 0 & 1 \\
X^3\cdot g(X): & 1 & 0 & 1 & 1 & 0 & 0 & 0 \\
\cline{2-8}
\text{resto}_1: & 0 & 1 & 1 & 0 & 1 & 0 & 1 \\
\end{array}
\qquad\Rightarrow\qquad q(X)\mathrel{+}=X^3
$$
 
\textbf{Paso 2.} Nuevo término líder: $X^5$, entonces $X^{5-3}=X^2$. Se resta
$X^2\cdot g(X)=X^5+X^3+X^2$:
$$
\begin{array}{r|c c c c c c c}
 & X^6 & X^5 & X^4 & X^3 & X^2 & X^1 & X^0 \\
\hline
\text{resto}_1: & 0 & 1 & 1 & 0 & 1 & 0 & 1 \\
X^2\cdot g(X): & 0 & 1 & 0 & 1 & 1 & 0 & 0 \\
\cline{2-8}
\text{resto}_2: & 0 & 0 & 1 & 1 & 0 & 0 & 1 \\
\end{array}
\qquad\Rightarrow\qquad q(X)\mathrel{+}=X^2
$$
 
\textbf{Paso 3.} Nuevo término líder: $X^4$, entonces $X^{4-3}=X$. Se resta
$X\cdot g(X)=X^4+X^2+X$:
$$
\begin{array}{r|c c c c c c c}
 & X^6 & X^5 & X^4 & X^3 & X^2 & X^1 & X^0 \\
\hline
\text{resto}_2: & 0 & 0 & 1 & 1 & 0 & 0 & 1 \\
X\cdot g(X): & 0 & 0 & 1 & 0 & 1 & 1 & 0 \\
\cline{2-8}
\text{resto}_3: & 0 & 0 & 0 & 1 & 1 & 1 & 1 \\
\end{array}
\qquad\Rightarrow\qquad q(X)\mathrel{+}=X
$$
 
\textbf{Paso 4.} Nuevo término líder: $X^3$, igual al grado de $g$, entonces $X^{3-3}=1$.
Se resta $1\cdot g(X)=X^3+X+1$:
$$
\begin{array}{r|c c c c c c c}
 & X^6 & X^5 & X^4 & X^3 & X^2 & X^1 & X^0 \\
\hline
\text{resto}_3: & 0 & 0 & 0 & 1 & 1 & 1 & 1 \\
1\cdot g(X): & 0 & 0 & 0 & 1 & 0 & 1 & 1 \\
\cline{2-8}
\text{resto}_4: & 0 & 0 & 0 & 0 & 1 & 0 & 0 \\
\end{array}
\qquad\Rightarrow\qquad q(X)\mathrel{+}=1
$$
 
El grado de $\text{resto}_4$ (que es $X^2$) ya es menor que el grado de $g(X)$ (que es 3), así
que la división termina acá:
$$q(X) = X^3+X^2+X+1, \qquad r(X) = X^2.$$
 
\textbf{b)} Verificación de $f(X) = q(X)\,g(X) + r(X)$:
\begin{align*}
q(X)\,g(X) &= (X^3+X^2+X+1)(1+X+X^3) \\
&= X^6+X^5+X^3+X^2+1 \quad\text{(desarrollando y cancelando términos repetidos)} \\
q(X)\,g(X)+r(X) &= (X^6+X^5+X^3+X^2+1) + X^2 = X^6+X^5+X^3+X^2+1 = f(X). \quad \ok
\end{align*}


\subsection*{Ejercicio 3 Polinomios irreducibles y primitivos}

Se usa $p(X) = 1+X^3+X^4$ como primitivo para construir $GF(2^4)$.

\textbf{a)} Evaluando en los dos elementos de $GF(2)$:
$$p(0) = 1, \qquad p(1) = 1+1+1 = 1.$$
Como $p(0)\neq 0$ y $p(1)\neq 0$, $p(X)$ no tiene raíces en $GF(2)$.

\textbf{b)} No tener raíces no alcanza para grado 4: un polinomio de
grado 4 sin raíces todavía podría factorizarse como producto de dos irreducibles de grado 2,
y el único irreducible de grado 2 es $X^2+X+1$. Se divide $p(X)$ por $X^2+X+1$:
$$p(X) \div (X^2+X+1) \;\Rightarrow\; \text{resto} = X \neq 0.$$
Como el resto no es nulo, $p(X)$ no es divisible por $X^2+X+1$. Al no tener raíces en $GF(2)$
y no ser divisible por el único irreducible de grado 2, $p(X)$ no admite ninguna
factorización posible en factores de grado menor: es irreducible. 

\textbf{c)} De $p(\alpha)=1+\alpha^3+\alpha^4=0$ se despeja
$$\alpha^4 = 1+\alpha^3,$$
\newpage
\begin{table}[h]
\centering
\renewcommand{\arraystretch}{1.15}
\begin{tabular}{c l c}
\toprule
Potencia & Polinomio & 4-tupla $(b_3b_2b_1b_0)$ \\
\midrule
$0$            & $0$                             & 0000 \\
$\alpha^{0}$   & $1$                             & 0001 \\
$\alpha^{1}$   & $\alpha$                        & 0010 \\
$\alpha^{2}$   & $\alpha^2$                      & 0100 \\
$\alpha^{3}$   & $\alpha^3$                      & 1000 \\
$\alpha^{4}$   & $1+\alpha^3$                    & 1001 \\
$\alpha^{5}$   & $1+\alpha+\alpha^3$             & 1011 \\
$\alpha^{6}$   & $1+\alpha+\alpha^2+\alpha^3$    & 1111 \\
$\alpha^{7}$   & $1+\alpha+\alpha^2$             & 0111 \\
$\alpha^{8}$   & $\alpha+\alpha^2+\alpha^3$      & 1110 \\
$\alpha^{9}$   & $1+\alpha^2$                    & 0101 \\
$\alpha^{10}$  & $\alpha+\alpha^3$               & 1010 \\
$\alpha^{11}$  & $1+\alpha^2+\alpha^3$           & 1101 \\
$\alpha^{12}$  & $1+\alpha$                      & 0011 \\
$\alpha^{13}$  & $\alpha+\alpha^2$               & 0110 \\
$\alpha^{14}$  & $\alpha^2+\alpha^3$             & 1100 \\
\bottomrule
\end{tabular}
\label{tab:gf16}
\end{table}

Se verifica además que $\alpha^{15}=\alpha\cdot\alpha^{14}=\alpha(\alpha^2+\alpha^3)=\alpha^3+\alpha^4=\alpha^3+1+\alpha^3=1$,
cerrando el ciclo de los $2^4-1=15$ pasos.

\subsection*{Ejercicio 4 Suma y producto sobre $GF(2^4)$}

Sean $f(X)=\alpha^3X^2+\alpha^7X+\alpha^2$ y $g(X)=\alpha^5X+\alpha^{10}$.

\textbf{a)} Se suma coeficiente a coeficiente:
\begin{align*}
\text{término en } X^0:&\; \alpha^2+\alpha^{10} = \alpha^2 + (\alpha+\alpha^3) = \alpha+\alpha^2+\alpha^3 = \alpha^8 \\
\text{término en } X^1:&\; \alpha^7+\alpha^5 = (1+\alpha+\alpha^2)+(1+\alpha+\alpha^3) = \alpha^2+\alpha^3 = \alpha^{14}
\end{align*}
$$f(X)+g(X) = \alpha^3X^2 + \alpha^{14}X + \alpha^8.$$

\textbf{b)} Para el producto conviene distribuir y usar que multiplicar potencias de
$\alpha$ es sumar exponentes módulo 15:
\begin{align*}
f(X)\cdot g(X) &= (\alpha^3X^2+\alpha^7X+\alpha^2)(\alpha^5X+\alpha^{10}) \\
&= \alpha^{3+5}X^3 + \alpha^{3+10}X^2 + \alpha^{7+5}X^2 + \alpha^{7+10}X + \alpha^{2+5}X + \alpha^{2+10} \\
&= \alpha^{8}X^3 + \alpha^{13}X^2 + \alpha^{12}X^2 + \alpha^{17\bmod 15}X + \alpha^{7}X + \alpha^{12} \\
&= \alpha^{8}X^3 + (\alpha^{13}+\alpha^{12})X^2 + (\alpha^{2}+\alpha^{7})X + \alpha^{12}.
\end{align*}
Sumando los coeficientes que se repiten en cada grado (pasando a polinomio):
$\alpha^{13}+\alpha^{12} = (\alpha^2+\alpha)+(1+\alpha) = 1+\alpha^2 = \alpha^9$, y
$\alpha^{2}+\alpha^{7} = \alpha^2+(1+\alpha+\alpha^2) = 1+\alpha = \alpha^{12}$. Entonces:
$$f(X)\cdot g(X) = \alpha^8X^3 + \alpha^9X^2 + \alpha^{12}X + \alpha^{12}.$$
\newpage
\subsection*{Ejercicio 5 División y potencia sobre $GF(2^4)$}

\textbf{a)} Dividir $f(X)=\alpha^3X^3+\alpha^9X^2+X+\alpha^6$ por $g(X)=X+\alpha^5$.

\begin{center}
\renewcommand{\arraystretch}{1.3}
\begin{tabular}{c c c}
\toprule
Coeficiente de $f$ & Cálculo & Resultado \\
\midrule
$\alpha^3$ (líder) & $b_0=\alpha^3$ & $\alpha^3$ \\
$\alpha^9$ & $b_1=\alpha^9+\alpha^5\cdot b_0=\alpha^9+\alpha^8$ & $\alpha^5$ \\
$1$ ($=\alpha^0$) & $b_2=\alpha^0+\alpha^5\cdot b_1=\alpha^0+\alpha^{10}$ & $\alpha^5$ \\
$\alpha^6$ & $r=\alpha^6+\alpha^5\cdot b_2=\alpha^6+\alpha^{10}$ & $\alpha^9$ \\
\bottomrule
\end{tabular}
\end{center}

(Cada paso: $\alpha^9+\alpha^8=(1+\alpha^2)+(\alpha+\alpha^2+\alpha^3)=1+\alpha+\alpha^3=\alpha^5$;
$\alpha^0+\alpha^{10}=1+(\alpha+\alpha^3)=\alpha^5$, $1+\alpha+\alpha^3=\alpha^5$ en la tabla;
$\alpha^6+\alpha^{10}=(1+\alpha+\alpha^2+\alpha^3)+(\alpha+\alpha^3)=1+\alpha^2=\alpha^9$.)

$$q(X) = \alpha^3X^2+\alpha^5X+\alpha^5, \qquad r(X)=\alpha^9.$$


\textbf{b)} $h(X)=\alpha^3X^2+\alpha^{11}X+\alpha^6$. Para ver con claridad de dónde sale cada
cosa, llamo $a=\alpha^3$, $b=\alpha^{11}$, $c=\alpha^6$ a los tres coeficientes, así
$h(X)=aX^2+bX+c$. Multiplicar $h(X)$ por sí mismo es multiplicar un trinomio por otro: hay
$3\times3=9$ productos cruzados, uno por cada par de términos (fila $\times$ columna):
 
\begin{center}
\renewcommand{\arraystretch}{1.4}
\begin{tabular}{c|ccc}
$\times$ & $aX^2$ & $bX$ & $c$ \\
\hline
$aX^2$ & $a^2X^4$ & $abX^3$ & $acX^2$ \\
$bX$   & $abX^3$  & $b^2X^2$& $bcX$   \\
$c$    & $acX^2$  & $bcX$   & $c^2$
\end{tabular}
\end{center}
 
Agrupando por grado (cada diagonal de la tabla cae en la misma potencia de $X$):
$$h(X)^2 = a^2X^4 \;+\; \underbrace{(ab+ab)}_{X^3}X^3 \;+\; \underbrace{(ac+b^2+ac)}_{X^2}X^2 \;+\; \underbrace{(bc+bc)}_{X^1}X \;+\; c^2.$$
 
Cada término que aparece \emph{dos veces} entonces se cancela. Solo 
sobreviven los términos de la diagonal ($a^2$, $b^2$, $c^2$) y el término $ac$, que aparece
\emph{tres} veces:
$$h(X)^2 = a^2X^4 + b^2X^2 + c^2.$$
 
Ahora sí, reemplazando por los exponentes (sumando exponentes módulo 15 en cada producto):
\begin{align*}
a^2 &= \alpha^3\cdot\alpha^3 = \alpha^{3+3}=\alpha^6 \\
b^2 &= \alpha^{11}\cdot\alpha^{11} = \alpha^{11+11}=\alpha^{22\bmod15}=\alpha^7 \\
c^2 &= \alpha^6\cdot\alpha^6 = \alpha^{6+6}=\alpha^{12}
\end{align*}
$$h(X)^2 = \alpha^6X^4 + \alpha^7X^2 + \alpha^{12}.$$

\textbf{c)} Elevando cada coeficiente de $h(X)$ al cuadrado y poniendolo en el doble del grado:
$$(\alpha^3)^2 X^{4} + (\alpha^{11})^2 X^{2} + (\alpha^6)^2 X^{0}
= \alpha^{6}X^4 + \alpha^{22\bmod 15}X^2 + \alpha^{12} = \alpha^6X^4+\alpha^7X^2+\alpha^{12}$$ que coincide con el resultado de (b). 

\subsection*{Ejercicio 6 Búsqueda de raíces}

Se buscan las raíces de $f(X)=X^2+\alpha^5X+\alpha^9$ sobre $GF(2^4)$. Como el cuerpo es
finito, alcanza con evaluar $f$ en los 16 elementos y quedarse con los que anulan el
polinomio. Evaluando:
$$f(\alpha^{11}) = \alpha^{22\bmod15}+\alpha^{5}\alpha^{11}+\alpha^9 = \alpha^7+\alpha^{16\bmod15}+\alpha^9 = \alpha^7+\alpha^1+\alpha^9.$$
Pasando a polinomio: $(1+\alpha+\alpha^2)+\alpha+(1+\alpha^2) = 0$. Efectivamente $\alpha^{11}$
es raíz. De manera análoga se verifica que $\alpha^{13}$ también lo es. Las raíces son
$$X_1=\alpha^{11}, \qquad X_2=\alpha^{13}.$$

\subsection*{Ejercicio 7 Ortogonalidad sobre $GF(2)$}

\textbf{a)} Producto interno de $u=(1,1,0,1)$ y $v=(0,1,1,1)$ en $V_4$ sobre $GF(2)$:
$$u\cdot v = (1\cdot0)+(1\cdot1)+(0\cdot1)+(1\cdot1) = 0+1+0+1 = 0.$$

\textbf{b)} Se busca un vector no nulo de $V_2$ ortogonal a sí mismo. Tomando $x=(1,1)$:
$$x\cdot x = 1\cdot1+1\cdot1 = 1+1 = 0,$$
y $x\neq(0,0)$. Sobre los números reales esto es imposible: $x\cdot x=\sum x_i^2$ es una suma
de cuadrados no negativos, que sólo puede anularse si $x=0$. Sobre $GF(2)$ esa implicación se
rompe porque $1+1=0$: un vector con una cantidad
\emph{par} de coordenadas no nulas puede ser ortogonal a sí mismo sin ser el vector nulo.

\subsection*{Ejercicio 8 Matriz generadora, forma sistemática y matriz de verificación}

Sea el código lineal $(6,3)$ generado por
$$G = \begin{pmatrix} 1&1&0&0&1&1 \\ 0&1&1&1&0&1 \\ 0&0&1&0&1&1 \end{pmatrix}.$$

\textbf{a)} Se lleva $G$ a la forma sistemática $[I_3\mid P]$ (Unidad 3, Sección 2.3) mediante
operaciones elementales de fila. 

\emph{Paso 1:} $F_1 \leftarrow F_1+F_2$:
$$F_1+F_2 = (1,1,0,0,1,1)+(0,1,1,1,0,1) = (1,0,1,1,1,0).$$

\emph{Paso 2:} $F_1 \leftarrow F_1+F_3$ y $F_2 \leftarrow F_2+F_3$:
\begin{align*}
F_1+F_3 &= (1,0,1,1,1,0)+(0,0,1,0,1,1) = (1,0,0,1,0,1) \\
F_2+F_3 &= (0,1,1,1,0,1)+(0,0,1,0,1,1) = (0,1,0,1,1,0)
\end{align*}


$$G_{sist} = \begin{pmatrix} 1&0&0&1&0&1 \\ 0&1&0&1&1&0 \\ 0&0&1&0&1&1 \end{pmatrix} = [I_3 \mid P], \qquad
P = \begin{pmatrix} 1&0&1 \\ 1&1&0 \\ 0&1&1 \end{pmatrix}.$$

\textbf{b)} A partir de $P$ hago $H=[P^T \mid I_3]$:
$$P^T = \begin{pmatrix} 1&1&0 \\ 0&1&1 \\ 1&0&1 \end{pmatrix}
\quad\Longrightarrow\quad
H = \begin{pmatrix} 1&1&0&1&0&0 \\ 0&1&1&0&1&0 \\ 1&0&1&0&0&1 \end{pmatrix}.$$

\textbf{c)} 
\begin{align*}
\text{fila}_0(G_{sist})\cdot\text{fila}_0(H) &= (1,0,0,1,0,1)\cdot(1,1,0,1,0,0) \\
&= 1{\cdot}1+0{\cdot}1+0{\cdot}0+1{\cdot}1+0{\cdot}0+1{\cdot}0 = 1+1 = 0 \\[4pt]
\text{fila}_1(G_{sist})\cdot\text{fila}_2(H) &= (0,1,0,1,1,0)\cdot(1,0,1,0,0,1) \\
&= 0{\cdot}1+1{\cdot}0+0{\cdot}1+1{\cdot}0+1{\cdot}0+0{\cdot}1 = 0
\end{align*}

\textbf{d)} Se codifica $u=(1,0,1)$ como $v=u\cdot G_{sist}$:
$$v = 1\cdot(1,0,0,1,0,1) + 0\cdot(0,1,0,1,1,0) + 1\cdot(0,0,1,0,1,1) = (1,0,1,1,1,0).$$
Verificación $v\cdot H^T=0$: como $v$ es combinación lineal de las filas de $G_{sist}$ y
$G_{sist}H^T=0$ por el punto (c), se vería que $vH^T=0$. Calculándolo:
$$v\cdot H^T = (0,\,0,\,0).$$

\textbf{e)} Linealidad de la codificación. Sean $u_1=(1,1,0)$ y $u_2=(1,0,0)$:
\begin{align*}
c_1 = u_1\cdot G_{sist} &= (1,1,0,0,1,1) + (0,1,0,1,1,0) = (1,1,0,0,1,1) \\
&\quad\text{(directo: } u_1=(1,1,0)\Rightarrow c_1=\text{fila}_0+\text{fila}_1=(1,1,0,0,1,1)) \\
c_2 = u_2\cdot G_{sist} &= \text{fila}_0 = (1,0,0,1,0,1)
\end{align*}
Con $u_3=u_1+u_2=(0,1,0)$, se codifica $u_3$ directamente ($u_3=\text{fila}_1$):
$$c_3 = u_3\cdot G_{sist} = (0,1,0,1,1,0).$$
Sumando $c_1+c_2 = (1,1,0,0,1,1)+(1,0,0,1,0,1) = (0,1,0,1,1,0) = c_3$. 

Como $v=uG_{sist}$ es una transformación lineal, $u_1G_{sist}+u_2G_{sist} =
(u_1+u_2)G_{sist}$ para cualquier par de mensajes.

\end{document}
